\chapter{Introduction to Multimedia Compression}

\section{Table of Contents}

\begin{itemize}
  \item Core Idea
  \item Key Concepts
  \item Theory and Formulas
  \item Visual Schemes
  \item Examples
\end{itemize}

\section{Core Idea}

Multimedia compression is needed because raw image, video, and audio data rates are far above practical storage and transmission limits. Effective codecs reduce rate by exploiting:

\begin{itemize}
  \item \textbf{Statistical redundancy}: correlation between nearby samples or frames.
  \item \textbf{Perceptual redundancy}: details humans cannot see or hear well.
  \item \textbf{Semantic redundancy}: task-irrelevant content, especially when the receiver is a machine.
\end{itemize}

Compression is always a design tradeoff among \textbf{rate}, \textbf{quality}, \textbf{complexity}, \textbf{delay}, and \textbf{robustness}.

\begin{calloutbox}{red!70!black}{Main compression principle}
A codec should spend bits where errors are visible or useful, and remove information that is predictable, imperceptible, or irrelevant to the final task.
\end{calloutbox}

\section{Key Concepts}

\subsection{Multimedia Representation}

\begin{center}\fitwidth{%
\begin{tabular}{|l|l|l|}
\hline
\textbf{Signal} & \textbf{Discrete representation} & \textbf{Key dependency} \\ \hline
Gray image & $N \times M$ samples $f_{n,m}$ & Spatial correlation \\ \hline
Color image & Three components, usually RGB or YCbCr & Luma/chroma sensitivity \\ \hline
Video & Sequence of color images over time & Spatial and temporal correlation \\ \hline
Audio & Time signal or frequency components & Hearing threshold and masking \\ \hline
\end{tabular}}
\end{center}

For an image scanned in raster order:

\[
k = (n-1)M + m, \qquad f_{n,m} = f_k
\]

For digital video:

\[
I : (n,m,T,c) \rightarrow x \in \{0,1,\ldots,2^b-1\}
\]

where $(n,m)$ is position, $T$ is frame index, $c$ is color component, and $b$ is bit depth.

\subsection{Human Visual System}

\begin{itemize}
  \item \textbf{Cones}: 6-7 million, concentrated near the fovea, responsible for color and high spatial resolution under good illumination.
  \item \textbf{Rods}: 75-150 million, more sensitive in low light, mainly detect intensity with lower resolution.
  \item Perceived brightness is approximately \textbf{logarithmic} with physical intensity.
  \item The eye adapts to illumination, so it cannot use the full physical dynamic range at once.
\end{itemize}

\begin{center}\includegraphics[width=0.833\linewidth]{ch01/hvs-receptors.png}\end{center}
Retinal receptors explain why compression can treat brightness, color, and fine details differently.

\begin{calloutbox}{red!70!black}{Contrast Sensitivity Function (CSF)}
The \textbf{CSF} gives sensitivity to luminance changes as a function of spatial frequency.

\[
\text{sensitivity} = \frac{1}{\text{minimum detectable contrast}}
\]

Sensitivity is highest around \textbf{2-5 cycles/degree} and lower at very low and very high spatial frequencies.

\textbf{Compression meaning:} high-frequency errors are often less visible, so high-frequency transform coefficients can be quantized more coarsely.
\end{calloutbox}

\subsection{Color Perception and Color Spaces}

Visible light lies roughly between \textbf{400 nm and 700 nm}. Cone sensitivities are uneven:

\begin{itemize}
  \item Red-sensitive cones: about \textbf{65\%}, peak near 575 nm.
  \item Green-sensitive cones: about \textbf{33\%}, peak near 535 nm.
  \item Blue-sensitive cones: about \textbf{2\%}, peak near 445 nm.
\end{itemize}

Color perception is \textbf{tristimulus}: a color corresponds to the 3D response vector of the three cone classes.

\begin{center}\fitwidth{%
\begin{tabular}{|l|l|l|}
\hline
\textbf{Space} & \textbf{Components} & \textbf{Compression relevance} \\ \hline
\textbf{RGB} & Red, Green, Blue & Device-oriented; three full-resolution channels \\ \hline
\textbf{HSV} & Hue, Saturation, Value & Perceptual description of color \\ \hline
\textbf{YCbCr} & Luminance Y, chrominance Cb/Cr & Separates brightness from color; enables chroma subsampling \\ \hline
\end{tabular}}
\end{center}

\begin{calloutbox}{red!70!black}{ITU-R BT.601 RGB to YCbCr}
\[
\begin{bmatrix}
Y\\
Cb\\
Cr
\end{bmatrix}
=
\begin{bmatrix}
0.299 & 0.587 & 0.114\\
-0.1687 & -0.3313 & 0.5\\
0.5 & -0.4187 & -0.0813
\end{bmatrix}
\begin{bmatrix}
R\\
G\\
B
\end{bmatrix}
+
\begin{bmatrix}
0\\
128\\
128
\end{bmatrix}
\]

The luminance weights follow visual sensitivity: green contributes most, blue least. The $+128$ offset centers chroma components in the 8-bit range.
\end{calloutbox}

\subsection{Sound Perception}

Audio coding uses a psychoacoustical model. A pure tone

\[
x(t)=a\sin(2\pi f_1 t)
\]

has power

\[
\sigma^2 = \frac{a^2}{2}
\]

and excites several nerve fibers, not only one.

For auditory filter $k$:

\[
H_k(f)=A_k(f)e^{j\phi_k(f)}
\]

The response to tone $f_1$ is:

\[
y_k(t)=aA_k(f_1)\sin(2\pi f_1t+\phi_k(f_1))
\]

The \textbf{spreading function} is:

\[
S_E(k)=A_k^2(f_1)
\]

\begin{calloutbox}{red!70!black}{Hearing threshold}
$S_a(f)$ is the minimum power for a tone at frequency $f$ to be audible. It is lowest around \textbf{1-4 kHz}, where human hearing is most sensitive.

\textbf{Compression meaning:} components below the threshold can be discarded or hidden under quantization noise.
\end{calloutbox}

\begin{calloutbox}{red!70!black}{Critical band}
Sinusoids close in frequency are integrated by hearing. For $N$ sinusoids near $f_1$:

\[
\sum_i \sigma_i^2 > S_a(f_1)
\]

The frequency interval where this energy integration holds is the \textbf{critical band}.
\end{calloutbox}

\subsection{Chroma Subsampling}

Notation $J:a:b$ describes chroma samples relative to a luma reference block:

\begin{itemize}
  \item $J$: horizontal luma reference size, usually 4.
  \item $a$: chroma samples in the first line.
  \item $b$: chroma samples in the second line.
\end{itemize}

\begin{center}\fitwidth{%
\begin{tabular}{|l|l|l|l|}
\hline
\textbf{Scheme} & \textbf{Horizontal chroma resolution} & \textbf{Vertical chroma resolution} & \textbf{Use} \\ \hline
\textbf{4:4:4} & Full & Full & High quality, no chroma reduction \\ \hline
\textbf{4:2:2} & Half & Full & Professional video \\ \hline
\textbf{4:2:0} & Half & Half & Most common image/video coding format \\ \hline
\textbf{4:1:1} & Quarter & Full & Older video systems \\ \hline
\end{tabular}}
\end{center}

For \textbf{YCbCr 4:2:0}, Y is full resolution, Cb and Cr are half resolution horizontally and vertically:

\[
\frac{4+1+1}{4+4+4}=\frac{6}{12}=0.5
\]

This gives \textbf{50\% data reduction} compared with full RGB/YUV 4:4:4 before any transform or entropy coding.

\subsection{Compression Tools}

\begin{center}\fitwidth{%
\begin{tabular}{|l|l|l|}
\hline
\textbf{Tool} & \textbf{Role} & \textbf{Lossy?} \\ \hline
\textbf{Transform} & Concentrates energy in few coefficients, e.g. DCT, wavelets, neural transforms & No by itself \\ \hline
\textbf{Prediction} & Removes spatial or temporal redundancy & No by itself \\ \hline
\textbf{Quantization} & Maps values to fewer levels; reduces precision & \textbf{Yes} \\ \hline
\textbf{Entropy coding} & Removes residual statistical redundancy, e.g. VLC, arithmetic coding & No \\ \hline
\end{tabular}}
\end{center}

\begin{calloutbox}{red!70!black}{Quantization is the lossy step}
In classical codecs, transform, prediction, and entropy coding can be reversible. The irreversible rate reduction comes from \textbf{quantization}.
\end{calloutbox}

\subsection{Lossless vs. Lossy Compression}

\begin{center}\fitwidth{%
\begin{tabular}{|l|l|l|}
\hline
\textbf{Property} & \textbf{Lossless} & \textbf{Lossy} \\ \hline
Reconstruction & Exact & Approximate \\ \hline
Exploits & Statistical redundancy & Statistical + perceptual redundancy \\ \hline
Main tool & Entropy coding, reversible prediction/transform & Quantization plus perceptual modeling \\ \hline
Typical image ratio & $T \leq 3$ & $T \approx 5$ or higher \\ \hline
Typical video ratio & Limited & $T \approx 20$ or higher \\ \hline
\end{tabular}}
\end{center}

\subsection{Machine-Centric Multimedia}

Not all multimedia is consumed by humans. Cameras in surveillance, IoT, industry, and robotics often feed algorithms.

\begin{center}\fitwidth{%
\begin{tabular}{|l|l|l|}
\hline
\textbf{Receiver} & \textbf{Priority} & \textbf{Compression goal} \\ \hline
Human & Perceptual quality & Hide distortion below HVS/HAS limits \\ \hline
Machine & Task accuracy & Preserve semantic features needed by model \\ \hline
\end{tabular}}
\end{center}

\textbf{Task-oriented communication} can discard visually important but task-irrelevant information if the downstream algorithm still performs well.

\section{Theory and Formulas}

\subsection{Quantization and Bit Depth}

If a scalar sample is represented with $L$ levels:

\[
b = \log_2 L
\]

where $b$ is bits per sample. Standard 8-bit components have:

\begin{itemize}
  \item $L=256$ gray levels.
  \item RGB: $256^3 \approx 16$ million colors.
  \item 24 bpp for full RGB.
\end{itemize}

HDR content may use 32-64 bits per channel.

\subsection{Rate and Compression Ratio}

\begin{calloutbox}{red!70!black}{Compression ratio}
\[
T = \frac{B_{\text{in}}}{B_{\text{out}}}
  = \frac{R_{\text{in}}}{R_{\text{out}}}
\]

$T$ measures how many times the coded representation is smaller than the original.
\end{calloutbox}

Coding rate:

\[
R_{\text{image}} = \frac{B_{\text{out}}}{NM} \quad [\text{bpp}]
\]

\[
R_{\text{audio/video}} = \frac{B_{\text{out}}}{T} \quad [\text{bps}]
\]

Typical values:

\begin{center}\fitwidth{%
\begin{tabular}{|l|l|}
\hline
\textbf{Application} & \textbf{Typical compression ratio} \\ \hline
Lossless image coding & $T \leq 3$ \\ \hline
Lossy image coding & $T \approx 5$ to much higher \\ \hline
Lossy video coding & $T \approx 20$ to much higher \\ \hline
\end{tabular}}
\end{center}

\subsection{Objective Quality Metrics}

Let $f$ be the original image and $\tilde{f}$ the decoded image. The error image is:

\[
\mathcal{E}(f,\tilde{f}) = f-\tilde{f}
\]

\begin{calloutbox}{red!70!black}{MSE and PSNR}
\[
\mathcal{D}(f,\tilde{f}) =
\frac{1}{NM}\|\mathcal{E}\|^2 =
\frac{1}{NM}\sum_{n,m}\mathcal{E}_{n,m}^2
\]

\[
\text{PSNR}(f,\tilde{f}) =
10\log_{10}\left(\frac{255^2}{\mathcal{D}(f,\tilde{f})}\right)
\quad [\text{dB}]
\]

\textbf{Meaning:} MSE/PSNR are simple and analytically useful, but they do not model perception.
\end{calloutbox}

Weighted PSNR introduces a filter $h$ that weights errors according to perceptual sensitivity:

\[
\mathcal{D}_W = \frac{1}{NM}\|h * \mathcal{E}\|^2
\]

\[
\text{WPSNR}(f,\tilde{f}) =
10\log_{10}\left(\frac{255^2}{\mathcal{D}_W}\right)
\]

\subsection{SSIM}

\begin{calloutbox}{red!70!black}{Structural Similarity Index}
For image blocks $x$ and $y$:

\[
\text{SSIM}(x,y)=
\frac{(2\mu_x\mu_y+C_1)(2\sigma_{xy}+C_2)}
{(\mu_x^2+\mu_y^2+C_1)(\sigma_x^2+\sigma_y^2+C_2)}
\]

SSIM combines:

\[
l(x,y)=\frac{2\mu_x\mu_y+C_1}{\mu_x^2+\mu_y^2+C_1}
\]

\[
c(x,y)=\frac{2\sigma_x\sigma_y+C_2}{\sigma_x^2+\sigma_y^2+C_2}
\]

\[
s(x,y)=\frac{\sigma_{xy}+C_3}{\sigma_x\sigma_y+C_3}
\]

Image SSIM is the average over blocks. Range is $[0,1]$, with 1 meaning identical.
\end{calloutbox}

\subsection{Learned Perceptual Metrics}

\textbf{LPIPS} (\textit{Learned Perceptual Image Patch Similarity}) compares images in a neural feature space, often using pre-trained VGG or ViT features. It is useful when pixel metrics fail, especially for highly compressed, generated, or neural-coded images.

\subsection{Subjective Quality}

Subjective tests with human observers are the reference for perceived quality, but they are expensive and statistically demanding. Objective metrics are proxies:

\begin{itemize}
  \item \textbf{PSNR}: easy, pixel-based.
  \item \textbf{SSIM}: structure-aware.
  \item \textbf{LPIPS}: learned perceptual feature distance.
\end{itemize}

\subsection{Complexity, Delay, and Robustness}

Beyond rate and quality, codecs must balance:

\begin{itemize}
  \item \textbf{Complexity}: real-time feasibility, hardware cost, power consumption.
  \item \textbf{Delay}: especially encoder delay; often depends on coding order and look-ahead.
  \item \textbf{Robustness}: sensitivity of compressed bitstreams to channel loss or packet errors.
\end{itemize}

\begin{center}\fitwidth{%
\begin{tabular}{|l|l|}
\hline
\textbf{Goal} & \textbf{Typical conflict} \\ \hline
Higher quality & Higher rate or complexity \\ \hline
Lower rate & Lower quality or higher delay \\ \hline
Lower delay & Less look-ahead, weaker compression \\ \hline
Higher robustness & More redundancy or lower compression efficiency \\ \hline
\end{tabular}}
\end{center}

\section{Visual Schemes}

\subsection{Visual and Auditory Perception}

\begin{center}\includegraphics[width=0.833\linewidth]{ch01/contrast-sensitivity-function.png}\end{center}
CSF shows why compression can quantize frequencies differently: mid spatial frequencies are most visible, while very low and very high frequencies can tolerate larger errors.

\begin{center}\includegraphics[width=0.833\linewidth]{ch01/hearing-threshold.png}\end{center}
The hearing threshold identifies inaudible components that can be removed or used to hide quantization noise.

\begin{center}\includegraphics[width=0.667\linewidth]{ch01/critical-band.png}\end{center}
Critical bands explain why nearby tones interact and why audio coders group frequency components.

\begin{center}\includegraphics[width=0.833\linewidth]{ch01/frequency-masking.png}\end{center}
Frequency masking allows a strong tone to hide weaker tones nearby in frequency.

\begin{center}\includegraphics[width=0.833\linewidth]{ch01/temporal-masking.png}\end{center}
Temporal masking permits distortion shortly before and after a strong sound, especially after the masker.

\subsection{Color and Chroma Sampling}

\begin{center}\includegraphics[width=\linewidth]{ch01/chroma-subsampling.png}\end{center}
Chroma subsampling reduces color resolution because humans are more sensitive to luminance than chrominance.

\subsection{Compression Pipeline}

\begin{center}\includegraphics[width=0.833\linewidth]{ch01/compression-tools.png}\end{center}
This block scheme summarizes the main codec tools: prediction/transform reduce redundancy, quantization reduces rate, and entropy coding packs symbols efficiently.

\begin{center}\footnotesize\itshape Diagram (Mermaid source):\end{center}
\begin{verbatim}
flowchart LR
    A[Input samples] --> B[Prediction or transform]
    B --> C[Quantization]
    C --> D[Entropy coding]
    D --> E[Bitstream]
    C --> F[Inverse quantization]
    F --> G[Reconstruction for prediction]
\end{verbatim}

\subsection{Quality Assessment Chain}

\begin{center}\includegraphics[width=0.833\linewidth]{ch01/weighted-psnr.png}\end{center}
Weighted PSNR filters the error before measuring energy, so distortions are penalized according to perceptual visibility.

\begin{center}\footnotesize\itshape Diagram (Mermaid source):\end{center}
\begin{verbatim}
flowchart LR
    A[Original image] --> C[Error image]
    B[Decoded image] --> C
    C --> D[Perceptual weighting]
    D --> E[MSE or WPSNR]
    C --> F[SSIM / LPIPS]
    F --> G[Perceptual quality estimate]
\end{verbatim}

\subsection{Video Perception}

\begin{center}\includegraphics[width=0.833\linewidth]{ch01/spatio-temporal-csf.png}\end{center}
Video perception depends on spatial and temporal frequencies; scene changes create masking and can temporarily hide distortion.

\section{Examples}

\begin{calloutbox}{green!55!black}{Uncompressed HD DVB data rate}
\textbf{Setup:} one luma component $1920 \times 1080$, two chroma components $960 \times 540$, 8 bits/sample, 50 fps.

\[
R = (1920\cdot1080 + 2\cdot960\cdot540)\cdot8\cdot50
\approx 1.24\ \text{Gbps}
\]

A 2-hour movie needs about \textbf{1.12 TB} uncompressed.

\textbf{Takeaway:} raw video is too large; compression is not optional.
\end{calloutbox}

\begin{calloutbox}{green!55!black}{4:2:0 data reduction}
Full 4:4:4 stores 4 Y samples, 4 Cb samples, and 4 Cr samples in a reference area. 4:2:0 stores 4 Y samples, 1 Cb sample, and 1 Cr sample:

\[
\frac{4+1+1}{4+4+4} = 0.5
\]

\textbf{Takeaway:} chroma subsampling alone halves the raw color data.
\end{calloutbox}

\begin{calloutbox}{green!55!black}{Same MSE, different perception}
Source examples compare several distortions with similar or equal MSE.

\begin{center}\fitwidth{%
\begin{tabular}{|l|l|l|}
\hline
\textbf{Error type} & \textbf{MSE} & \textbf{SSIM} \\ \hline
Distributed white noise, $\sigma=4$ & 16 & 0.906 \\ \hline
Concentrated on $100\times100$ pixels & 16 & 0.972 \\ \hline
Concentrated on contours & 16 & 0.987 \\ \hline
Noise on high spatial frequencies & 16 & 0.882 \\ \hline
Chroma subsampling & 21.27 & - \\ \hline
\end{tabular}}
\end{center}

\textbf{Takeaway:} MSE alone is misleading because it ignores where and how errors appear.
\end{calloutbox}

\begin{calloutbox}{green!55!black}{Codec comparison at about 50 KB}
The source compares codecs on the same image target size.

\begin{center}\fitwidth{%
\begin{tabular}{|l|l|l|l|l|}
\hline
\textbf{Codec} & \textbf{Year} & \textbf{Size (Bytes)} & \textbf{Ratio} & \textbf{PSNR} \\ \hline
Original & - & 6,220,800 & 1.0:1 & - \\ \hline
PNG & 1996 & 2,923,960 & 2.1:1 & - \\ \hline
4:2:0 only & - & 3,110,400 & 2.0:1 & 43.05 dB \\ \hline
JPEG & 1991 & 49,795 & 124.9:1 & 26.89 dB \\ \hline
JPEG 2000 & - & 50,943 & 122.1:1 & 29.82 dB \\ \hline
WebP & 2010 & 42,676 & 145.8:1 & 30.88 dB \\ \hline
AVIF & 2019 & 48,745 & 127.6:1 & 32.45 dB \\ \hline
\end{tabular}}
\end{center}

\textbf{Takeaway:} at similar size, newer codecs can provide better objective quality.
\end{calloutbox}
